%
%  chapter2.tex
%
\chapter{State of the Art}
\label{chap:stateofart}

\section{Gamification in Education}
\label{sec:gamification}

Gamification is defined as the application of game design elements in non-game contexts \cite{deterding2011gamification}. In educational settings, these elements typically include points, leaderboards, badges, progress bars, and challenge structures. The underlying motivation is to harness the intrinsic drivers of game engagement---competition, mastery, and reward---to sustain learner effort in contexts where motivation often wanes.

\subsection{Evidence of Effectiveness}

Empirical evidence on gamification's effectiveness in education has accumulated steadily. A systematic literature review of 24 empirical studies found positive effects in the majority of cases, with social comparison mechanisms such as leaderboards and peer challenges producing the most consistent engagement gains \cite{hamari2014does}. However, the same review noted important moderating factors: the type of gamification element, the duration of the intervention, and the demographic profile of learners all influence outcomes. These nuances informed the design of U!REKA Play4Change's scoring engine, which combines individual progression with configurable penalty and reward thresholds rather than relying solely on competitive ranking.

\section{Serious Games}
\label{sec:seriousgames}

Serious games are digital games designed with a primary purpose beyond pure entertainment, typically targeting learning, training, or behaviour change. Unlike simple gamification, serious games situate the learner inside a coherent game world with narrative context, thus providing richer motivational scaffolding \cite{zhonggen2019meta}.

A meta-analysis of serious game research spanning a decade found statistically significant improvements in learning outcomes across multiple disciplines, including science, health, and civic education \cite{zhonggen2019meta}. The analysis identified immediate feedback and clear goal structures as the two strongest predictors of learning gain. U!REKA Play4Change draws on these findings by structuring each daily task as a self-contained challenge with instant scoring feedback and a defined daily goal.

\section{Multiplatform Mobile Development}
\label{sec:multiplatform}

\subsection{Native versus Cross-Platform Approaches}

Mobile application development historically bifurcated into native development---separate codebases per platform using Swift or Objective-C for iOS and Java or Kotlin for Android---and cross-platform development, where a single codebase targets multiple platforms through an abstraction layer. Native development maximises access to platform APIs and delivers the highest performance, but doubles the maintenance burden for teams targeting both major mobile platforms. Cross-platform frameworks trade some low-level control for significantly reduced duplication.

Established cross-platform frameworks include React Native, Flutter, and Kotlin Multiplatform. React Native renders components via a JavaScript bridge, which introduces latency for complex interactions. Flutter compiles to native ARM code and provides its own rendering engine, achieving near-native performance but diverging visually from platform conventions. Kotlin Multiplatform takes a different approach: it shares only business logic across platforms while delegating the user interface to platform-native or declarative frameworks, enabling a balance between code sharing and native look and feel.

\subsection{Justification for Kotlin Multiplatform and Compose Multiplatform}

U!REKA Play4Change adopts Kotlin Multiplatform for shared business logic---including the scoring engine, data models, network clients, and local caching---combined with Compose Multiplatform for the shared declarative user interface. This combination allows Android and iOS to share both logic and UI components from a single codebase while maintaining access to native APIs where required. The choice reduces duplication and accelerates feature delivery without sacrificing the performance characteristics required for smooth animated game interfaces.

\section{AI-Generated Content in Education}
\label{sec:aicontent}

\subsection{Large Language Models for Educational Content}

Large language models (LLMs) have demonstrated strong capability in generating coherent, contextually appropriate text across a wide range of domains. Their application to educational content generation is a rapidly growing area: LLMs can produce question-and-answer pairs, explanatory summaries, and scenario-based tasks at scale and at low marginal cost compared to human authoring. The primary challenges are factual accuracy, content diversity, and deduplication, all of which require careful pipeline design.

\subsection{LangChain4j and the Mistral API}

LangChain4j is a Java library that provides abstractions for building LLM-powered applications, including prompt templating, model chaining, and output parsing. It supports multiple model backends, including the Mistral API, which provides instruction-tuned models well suited to structured task generation \cite{langchain4j2024, mistral2024}. U!REKA Play4Change uses LangChain4j to orchestrate nightly batch jobs that generate daily tasks parameterised by subject domain and difficulty level. Generated tasks are embedded using a sentence-level embedding model and compared against existing stored embeddings via pgvector to detect and discard near-duplicate content before persistence \cite{pgvector2024}.

\section{Related Platforms}
\label{sec:relatedplatforms}

Several gamified learning platforms have gained widespread adoption, including Duolingo, Kahoot, and Quizlet. Duolingo applies a streak-based engagement model to language learning and demonstrates strong retention at scale; however, its content is entirely hand-authored by domain specialists, limiting expansion to new topics. Kahoot focuses on synchronous, instructor-led quiz sessions and is therefore unsuitable for self-paced asynchronous learning. Quizlet allows community-generated flashcard decks but does not provide an adaptive scoring engine or AI-generated content.

None of these platforms combines multiplatform delivery from a single codebase with automated AI-driven content generation. U!REKA Play4Change addresses this gap by integrating both capabilities, enabling educational networks to deploy a fully configured, self-sustaining gamified platform for any subject domain without sustained content authoring effort.
